% Conclusion. Part of the team paper; included by main.tex.
\section{Discussion and conclusion}
\label{sec:conclusion}

We moved the in-context learning setup of Garg et al.\ from function regression to closed-loop control, and found that a small
transformer infers a hidden goal from demonstrations well enough to match goal-seeing teachers on new goals from the training
distribution.

\textbf{What we learn.} The context acts as a strong control channel: the model follows demonstrations of the wrong goal as
confidently as the right ones. Offline error alone cannot certify a controller, so closed-loop evaluation is necessary. Out of
distribution, coarse behavior (reaching the goal) transfers before precise behavior (holding the block there).

\textbf{Did our hypotheses hold?} %CHECK: keep in sync with the hypotheses in Section 1
(i) Yes: in distribution the transformer matches the teachers, and the controls fall to 0--2\%.
(ii) Partly: reaching transfers to the held-out quadrant, but success-at-end drops by about 20--25 points, more than we expected.
(iii) No: success does not keep rising with $k$ but saturates by $k = 1$--$2$. This fits a goal that one successful demonstration
nearly reveals.
(iv) Not reliably: a pilot that already preferred correct demonstrations offline scored 0\% in closed loop. Offline error tracked
closed-loop success only once it was small.

\textbf{Limitations.} The goal is a 2-D position that a successful demonstration nearly reveals in its final state, so the
inference required is simple. We are therefore building a baseline that estimates the goal from the final T pose in the
demonstrations and gives it to the BC teacher, which will show how much more the transformer does. The main results come from one
training seed, and an independent run shows that success-at-end on held-out goals varies widely across seeds, so we are training two
more seeds to obtain between-seed error bars. Finally, v1 trains a single learner on a mixture of PPO and BC prompts, and because the
BC teacher is distilled from PPO, the two teachers implement nearly the same policy. The held-out quadrant tests
recombination more than extrapolation: each coordinate's sign appears in $Q_1$--$Q_3$, and only the combination is new. Every
demonstration and query state comes from clean teacher rollouts, so the learner never trains on the states it drifts into; Smith
et al.~\cite{smith2026transformers} found that removing state perturbations from training collapsed in-context control, which
suggests this as one cause of the holding failure on $Q_4$.

\textbf{Next steps.} v2 will treat PPO and BC policies as distinct teacher types, analogous to how Garg et al.\ train one model per
function class. We will retrain the full-state BC teacher on demonstration data in simulation instead of distilling it from PPO, so that its
behavior no longer derives from the PPO teacher. We will then train a separate transformer on each class, with the architecture,
data sizes and schedule of v1, so that only the training class changes, and evaluate each learner on its own class, on the other
class, and next to the v1 mixed learner. Three further experiments target the open questions above. Training on query states from
noisy teacher rollouts with clean labels addresses the holding failure on held-out goals. A sweep over the number of training goals
will locate where in-context generalization emerges~\cite{raventos2023pretraining}. Hidden variables that the final state does not
reveal, such as the T's mass or friction, will test inference beyond reading off the endpoint. Finally, we will repeat the study
with the camera-observation teachers.
